\begin{abstract}
\noindent Робот, который объезжает объекты по открытословарной семантической карте, получает две ошибки сразу. Детектор неточно размечает объекты: недорисовывает их или путает класс. Оценка позы дрейфует, и объекты попадают на карту не туда. Известные статистические гарантии для планирования по картам восприятия учитывают только первую ошибку и считают позу робота известной.

Мы калибруем одну величину --- расстояние промаха: насколько истинный объект выходит за свою область на карте, построенной по оценённым позам. Её конформный квантиль по сценам задаёт запас. Любой путь, который держится от области класса на безопасном расстоянии плюс этот запас, безопасен с заданной вероятностью, какой бы планировщик его ни построил.

На 21 сцене бенчмарка OSMa-Bench совместный запас держит заданное покрытие при любом уровне дрейфа. Запасы, откалиброванные только по меткам или только по позе, его теряют. В условиях ближайшей работы (известная поза, закрытый словарь) метод не уступает калибровке меток, а при дрейфе и открытом словаре превосходит её. Цену гарантии задаёт наихудшая ошибка восприятия. С реальной визуальной одометрией без замыканий циклов метод отказывается давать гарантию, вместо того чтобы выдать ложную.
\end{abstract}
